\subsection{EXTENSION OF DERIVATIONS}

Lemma 4. Let V be a K-module and T the tensor algebra of V. Let u be an endomorphism of V. There exists one and only one derivation of T which extends u. This derivation commutes with the symmetry operators on T.

Let \(F = V \times V \times \cdots \times V\) (n factors). The mapping

\[
\begin{array}{r l} (x _ {1}, \ldots , x _ {n}) & \mapsto u x _ {1} \otimes x _ {2} \otimes \dots \otimes x _ {n} \\ & \qquad + x _ {1} \otimes u x _ {2} \otimes \dots \otimes x _ {n} + \dots + x _ {1} \otimes x _ {2} \otimes \dots \otimes u x _ {n} \end{array}
\]

of F into \(\bigotimes^{n}\) V is multilinear. Hence there exists an endomorphism \(u_{n}\) of \(\bigotimes^{n}\) V such that:

\[
u _ {n} \left(x _ {1} \otimes \dots \otimes x _ {n}\right) = u x _ {1} \otimes \dots \otimes x _ {n} + \dots + x _ {1} \otimes \dots \otimes u x _ {n}
\]

for all \(x_1, \ldots, x_n\) in V. Then \(u_1 = u\) . Let \(v\) be the endomorphism of the K-module T which coincides with \(u_n\) on each \(\mathrm{T}^n = \bigotimes^n \mathrm{V}\) and which is zero on \(\mathrm{T}^0 = \mathrm{K}. 1\) . We show that \(v\) is a derivation of T. If \(x_1, \ldots, x_n, y_1, \ldots, y_p\) are elements of V, then

\[
\begin{array}{l} v ((x _ {1} \otimes \dots \otimes x _ {n}) \otimes (y _ {1} \otimes \dots \otimes y _ {p})) \\ = \sum_ {i = 1} ^ {n} x _ {1} \otimes \dots \otimes x _ {i - 1} \otimes u x _ {i} \otimes x _ {i + 1} \otimes \dots \otimes x _ {n} \otimes y _ {1} \otimes \dots \otimes y _ {p} \\ \qquad + \sum_ {j = 1} ^ {p} x _ {1} \otimes \dots \otimes x _ {n} \otimes y _ {1} \otimes \dots \otimes y _ {j - 1} \otimes u y _ {j} \otimes y _ {j + 1} \otimes \dots \otimes y _ {p} \\ = v (x _ {1} \otimes \dots \otimes x _ {n}) \otimes (y _ {1} \otimes \dots \otimes y _ {p}) + (x _ {1} \otimes \dots \otimes x _ {n}) \otimes v (y _ {1} \otimes \dots \otimes y _ {p}). \end{array}
\]

By linearity it follows that \(v\) is a derivation. The uniqueness of \(v\) is obvious.

Finally, clearly \(u_{n}\) commutes with all the symmetry operators on \(\otimes V\) , whence the last assertion.

PROPOSITION 7. Let \(\mathfrak{g}\) be a Lie algebra, U its enveloping algebra, \(\sigma\) the canonical mapping of \(\mathfrak{g}\) into U and D a derivation of \(\mathfrak{g}\) .

\begin{enumerate}
\def\labelenumi{(\alph{enumi})}
\item
  There exists one and only one derivation \(\mathrm{D_U}\) of U such that \(\sigma \circ \mathrm{D} = \mathrm{D_U} \circ \sigma\) (that is such that \(\mathrm{D_U}\) extends D, when \(\mathfrak{g}\) can be identified with a submodule of U under \(\sigma\) ).
\item
  \(D_{U}\) leaves stable \(U_{n}\) and the set \(U^{n}\) of images in U of the homogeneous symmetric tensors of order n over g.
\item
  \(\mathbf{D}_{\mathbf{U}}\) commutes with the principal antiautomorphism of \(\mathbf{U}\) .
\item
  If \(\mathbf{D}\) is the inner derivation of \(\mathfrak{g}\) defined by an element \(x\) of \(\mathfrak{g}\) , \(\mathrm{D}_{\mathbf{U}}\) is the inner derivation of \(\mathbf{U}\) defined by \(\sigma(x)\) .
\end{enumerate}

Let \(D_T\) be the derivation of the tensor algebra \(T\) of \(g\) which extends \(D\) (Lemma 4). The two-sided ideal \(J\) of \(T\) generated by the

\[
x \otimes y - y \otimes x - [ x, y ]
\]

\((x,y \text{ in } \mathfrak{g})\) is stable under \(\mathrm{D}_{\mathrm{T}}\) . For:

\[
\begin{array}{r} \mathrm{D} _ {\mathrm{T}} (x \otimes y - y \otimes x - [ x, y ]) = \mathrm{D} x \otimes y - y \otimes \mathrm{D} x - [ \mathrm{D} x, y ] \\ + x \otimes \mathrm{D} y - \mathrm{D} y \otimes x - [ x, \mathrm{D} y ]. \end{array}
\]

On passing to the quotient, \(D_{T}\) defines a derivation \(D_{U}\) of U such that \(\sigma \circ D = D_{U} \circ \sigma\) . The uniqueness of \(D_{U}\) is immediate since 1 and \(\sigma(\mathfrak{g})\) generate the algebra U. Assertion (b) is obvious. Let A be the principal antiautomorphism of U. We now prove (c). If \(x_{1}, \ldots, x_{n}\) are in g, then

\[
\begin{array}{r l} \mathrm{D} _ {\mathrm{U}} \mathrm{A} (\sigma (x _ {1}) \dots \sigma (x _ {n})) & = \mathrm{D} _ {\mathrm{U}} ((- 1) ^ {n} \sigma (x _ {n}) \dots \sigma (x _ {1})) \\ & = (- 1) ^ {n} \sum_ {i = 1} ^ {n} \sigma (x _ {n}) \dots \mathrm{D} _ {\mathrm{U}} (\sigma (x _ {i})) \dots \sigma (x _ {1}) \\ & = (- 1) ^ {n} \sum_ {i = 1} ^ {n} \sigma (x _ {n}) \dots \sigma (\mathrm{D} x _ {i}) \dots \sigma (x _ {1}) \\ & = \mathrm{A} \bigg (\sum_ {i = 1} ^ {n} \sigma (x _ {1}) \dots \sigma (\mathrm{D} x _ {i}) \dots \sigma (x _ {n}) \bigg) \\ & = \mathrm{AD} _ {\mathrm{U}} (\sigma (x _ {1}) \dots \sigma (x _ {n})). \end{array}
\]

Finally, let \(x \in \mathfrak{g}\) . Let \(\Delta\) be the inner derivation \(y \mapsto \sigma(x)y - y\sigma(x)\) of U (Algebra, Chapter IV, § 4, no. 3, Example 2). Then, for \(x' \in \mathfrak{g}\) ,

\[
(\Delta \circ \sigma) (x ^ {\prime}) = \sigma (x) \sigma (x ^ {\prime}) - \sigma (x ^ {\prime}) \sigma (x) = \sigma ([ x, x ^ {\prime} ]) = (\sigma \circ \operatorname{ad} x) (x ^ {\prime}),
\]

whence \(\Delta \circ \sigma = \sigma \circ \operatorname{ad} x\) . This completes the proof.

Applying Proposition 7 to the case of a commutative Lie algebra, it is seen that every endomorphism u of a K-module can be extended uniquely to a derivation of the symmetric algebra of this module; this derivation is derived on passing to the quotient from the derivation of the tensor algebra which extends u.

We again take a Lie algebra g over K and let D be a derivation of g. We use the earlier notation T, S, U, G. Let \(D_{T}\) , \(D_{S}\) be the derivations of T, S which extend D and let \(D_{U}\) be the unique derivation of U such that \(\sigma \circ D = D_{U} \circ \sigma\) . Since \(D_{U}\) leaves the \(U_{n}\) stable, \(D_{U}\) defines on taking quotients a derivation \(D_{G}\) of G. Since \(D_{U}\) and \(D_{S}\) are derived from \(D_{T}\) when passing to quotients, the commutative diagram (3) proves that \(D_{G}\) can also be derived from \(D_{S}\) by the homomorphism \(\omega\) defined in no. 6. If further K is a field of characteristic 0, the isomorphisms of diagram (4) map one into another the restrictions of \(D_{T}\) , \(D_{S}\) , \(D_{U}\) , \(D_{G}\) to \(S^{n}\) , \(S^{n}\) , \(U^{n}\) , \(G^{n}\) . Hence the canonical isomorphism of S onto U maps \(D_{S}\) to \(D_{U}\) .

